\documentclass[10pt,a4paper]{amsart}
\usepackage[margin=19mm]{geometry}
\usepackage{amsmath,amssymb,mathtools}
\usepackage[scr=rsfs,cal=euler]{mathalfa}
\usepackage{xfrac}
\usepackage[unicode,hidelinks,bookmarks=false]{hyperref}
\newcommand{\ie}{i.e.\ }
\newcommand{\cf}{cf.\ }
\newcommand{\resp}{respectively\ }
\newcommand{\Subsection}{\subsection}
\providecommand{\nobreakdash}{\nobreak\hbox{-}}
\setlength{\emergencystretch}{2em}
\multlinegap0pt
\usepackage{amssymb}
\usepackage[shortlabels]{enumitem}
\newtheorem{lemma}{Lemma}
\newtheorem{theorem}{Theorem}
\newcommand{\med}{\operatorname{med}}
\title{List-distance consistent vertices in trees are confined to a path}
\author{Fei-Huang Chang, Ma-Lian Chia, David Kuo, Guan-Ting Lai}
\date{Source: arXiv:2609.03803v1 (CC BY 4.0)}
\begin{document}
\maketitle

\noindent\emph{Source and licence.} List-distance consistent vertices in trees are confined to a path. Version arXiv:2609.03803v1 (CC BY 4.0), distributed by arXiv under the Creative Commons Attribution 4.0 International licence, http://creativecommons.org/licenses/by/4.0/, as recorded in the arXiv licence metadata for this exact version and shown on the abstract page https://arxiv.org/abs/2609.03803v1. Full licence notice: \texttt{licenses/arxiv-2609.03803-CC-BY-4.0.txt}.\par\smallskip

\noindent\textit{Editorial scope.} Counted unit: the article's theorem thm:path (source0.tex lines 345--360). The article's complete original proof of it is delivered with it (source0.tex lines 362--432, one \texttt{proof} environment of 71 lines). No mathematics is written, repaired or re-derived: the statement and the proof are the article's own lines, in the article's own notation, and no retained line is substituted or re-flowed.  Retained uncounted as the source-local context the proof needs: lines 284--292; lines 323--328.  Nothing was omitted from any retained span, so the package carries no omission record.  No retained line contains a citation, so the wrapper carries no bibliography and no bibliography entry.  No retained line contains a figure environment, an image-inclusion command or drawing code, so no image is delivered and the package declares no asset dependency.  Editorial formatting only, in addition: an AMS \texttt{amsart} preamble that restates the article's own declarations and macros used by the retained lines, namely the environments \texttt{lemma}, \texttt{theorem} on separate counters, together with the standard packages those lines need.  The retained environments print in this document as Lemma 1, Theorem 1 in order of appearance, whereas the article prints the same items with the numbering of its own full document.  The complete native source of the article as distributed in the arXiv e-print for this exact version is bundled unchanged as \texttt{sources/209316/source0.tex}.
\begin{lemma}
\label{lem:monotone}
Let $c$ be a labeling of a connected graph $G$ and let $u,v\in S_c(G)$ with
$c(u)<c(v)$. If $P:u=u_1u_2\cdots u_k=v$ is a shortest $u$--$v$ path in $G$,
then
\[
  c(u_1)<c(u_2)<\cdots<c(u_k).
\]
\end{lemma}

\begin{lemma}
\label{lem:threebranch}
Let $G$ be a connected graph and let $u_1,u_2,u_3$ lie in three distinct
components of $G-w$ for some vertex $w$. Then for every labeling $c$ of $G$ at
least one of $u_1,u_2,u_3$ fails to be consistent under $c$.
\end{lemma}

\begin{theorem}
\label{thm:path}
Let $T$ be a tree, let $c$ be a labeling of $T$, and put $S=S_c(T)$. Assume
$|S|\ge2$, and let $a,b\in S$ be the vertices with
\[
  c(a)=\min\{c(x):x\in S\},\qquad c(b)=\max\{c(x):x\in S\}.
\]
Write $P_{ab}=p_1p_2\cdots p_m$ with $p_1=a$ and $p_m=b$. Then
\begin{enumerate}
\item[\textup{(i)}] $S\subseteq V(P_{ab})$;
\item[\textup{(ii)}] $c(p_i)=c(a)+i-1$ for $1\le i\le m$; in particular
$c(b)=c(a)+m-1$ and $c\bigl(V(P_{ab})\bigr)=[\,c(a),c(b)\,]$;
\item[\textup{(iii)}] every $z\in V(T)\setminus V(P_{ab})$ satisfies
$c(z)<c(a)$ or $c(z)>c(b)$.
\end{enumerate}
\end{theorem}

\begin{proof}
\textbf{Step 1: proof of (i).}
Suppose there is $x\in S\setminus V(P_{ab})$, and let $w=\med(a,b,x)$. Since
$w\in V(P_{ab})$ and $x\notin V(P_{ab})$ we have $w\ne x$. We distinguish
three cases.

If $w\notin\{a,b\}$, then $w$ is an interior vertex of $P_{ab}$, and $a$, $b$,
$x$ lie in three distinct components of $T-w$. As $a,b,x\in S$, this
contradicts Lemma~\ref{lem:threebranch}.

If $w=a$, then $a$ lies on $P_{xb}$ and $a\notin\{x,b\}$, so $a$ is an interior
vertex of $P_{xb}$. Since $x,b\in S$, Lemma~\ref{lem:monotone} applied to $x$
and $b$ shows that $c(a)$ lies strictly between $c(x)$ and $c(b)$. This
contradicts the minimality of $c(a)$ on $S$, because $x,b\in S$ forces
$c(a)<c(x)$ and $c(a)<c(b)$.

If $w=b$, the symmetric argument shows that $c(b)$ lies strictly between $c(x)$
and $c(a)$, contradicting the maximality of $c(b)$ on $S$. Hence
$S\subseteq V(P_{ab})$.

\medskip
\noindent\textbf{Step 2: the labels increase along $P_{ab}$.}
The path $P_{ab}$ is the unique, hence shortest, $a$--$b$ path in $T$, and
$a,b\in S$ with $c(a)<c(b)$. Lemma~\ref{lem:monotone} therefore gives
\begin{equation}
  c(a)=c(p_1)<c(p_2)<\cdots<c(p_m)=c(b),
  \label{eq:increase}
\end{equation}
and in particular $c(a)\le c(p_i)\le c(b)$ for all $i$.

\medskip
\noindent\textbf{Step 3: proof of (iii).}
Suppose $z\notin V(P_{ab})$ satisfies $c(a)<c(z)<c(b)$, and let
$p_i=\med(a,b,z)$, the vertex of $P_{ab}$ nearest to $z$; put
$r=d(p_i,z)\ge1$. Every $a$--$z$ path passes through $p_i$, and likewise for
$b$, so
\[
  d(a,z)=(i-1)+r>i-1=d(a,p_i),
  \qquad
  d(b,z)=(m-i)+r>m-i=d(b,p_i).
\]

Consistency of $a$ applied to the pair $(p_i,z)$ gives $c(a,p_i)\le c(a,z)$. By
\eqref{eq:increase} we have $c(p_i)\ge c(a)$, and by assumption $c(z)>c(a)$, so
both absolute values may be resolved and we obtain
\[
  c(p_i)-c(a)\le c(z)-c(a),
  \qquad\text{that is,}\qquad c(p_i)\le c(z).
\]

Consistency of $b$ applied to the same pair gives $c(b,p_i)\le c(b,z)$. By
\eqref{eq:increase} we have $c(p_i)\le c(b)$, and by assumption $c(z)<c(b)$,
whence
\[
  c(b)-c(p_i)\le c(b)-c(z),
  \qquad\text{that is,}\qquad c(p_i)\ge c(z).
\]

Combining the two inequalities just obtained yields $c(p_i)=c(z)$. But
$p_i\in V(P_{ab})$ and $z\notin V(P_{ab})$, so $p_i\ne z$, contradicting the
injectivity of $c$. This proves (iii).

\medskip
\noindent\textbf{Step 4: proof of (ii).}
By Step 3, every vertex whose label lies in the interval $[\,c(a),c(b)\,]$
belongs to $V(P_{ab})$; hence $[\,c(a),c(b)\,]\subseteq c\bigl(V(P_{ab})\bigr)$.
By \eqref{eq:increase} the reverse inclusion also holds, so
$c\bigl(V(P_{ab})\bigr)=[\,c(a),c(b)\,]$. Comparing cardinalities gives
$c(b)-c(a)+1=m$, and \eqref{eq:increase} then forces $c(p_i)=c(a)+i-1$ for
every $i$.
\end{proof}

\end{document}
